%% ===================================================================
\section{Explicit form of BFKL\texorpdfstring{$\otimes$}{x}LO and JIMWLK\texorpdfstring{$\otimes$}{x}LO}\label{app:explicit}
%% ===================================================================

We write the LO spectrum~\eqref{eq:LO} in coordinate space as
\begin{align}
\frac{dN_{\rm LO}}{d^2w\,d^2w'\,dy}={}&\frac{\alpha_s}{4\pi^4(N_c^2-1)}\int_{x,y}\KK(x-w)\cdot\KK(y-w')\,\langle\rho^f(x)\rho^f(y)\rangle\notag\\
&\times\Tr\Big\{\big(\Ud(x)-\Ud(w)\big)\big(U(y)-U(w')\big)\Big\},
\label{eq:LOx}
\end{align}
so that $\frac{dN}{dy\,d^2\mathbf k}=k^+\frac{d^3N}{d^3k}=\int_{w,w'}e^{-ik\cdot(w'-w)}\frac{dN}{d^2w\,d^2w'\,dy}$. Here $y$ is a transverse position and not the rapidity, which appears only in $dy$.

\subsection{BFKL}

The BFKL equation for the charge correlator in coordinate space reads
\begin{align}
\frac{d\langle\rho^f(x)\rho^f(y)\rangle}{dY}=-\frac{\alpha_sN_c}{2\pi^2}\Bigg[&\int_zK_D(x,y,z)\,\langle\rho^f(x)\rho^f(y)\rangle-2\delta^{(2)}(x-y)\int_{z,z'}K(z,z',x)\,\langle\rho^f(z)\rho^f(z')\rangle\notag\\
&+\int_z\big(2K(y,z,x)-K(z,z,x)\big)\langle\rho^f(y)\rho^f(z)\rangle\notag\\
&+\int_z\big(2K(x,z,y)-K(z,z,y)\big)\langle\rho^f(x)\rho^f(z)\rangle\Bigg],
\label{eq:vBFKL}
\end{align}
with $K$ and $K_D$ of Eqs.~\eqref{eq:NX} and~\eqref{eq:KD}. Applied to Eq.~\eqref{eq:LOx}, after renaming integration variables to a common correlator, it gives
\begin{align}
\big[{\rm BFKL}\otimes{\rm LO}\big]\,\delta Y={}&-\frac{\alpha_s^2N_c\,\delta Y}{8\pi^6(N_c^2-1)}\int_{x,y,z}\langle\rho^f(x)\rho^f(y)\rangle\notag\\
&\times\Big[\KK(x-w)\cdot\KK(y-w')\,K_D(x,y,z)\,\Tr\big\{(\Ud(x)-\Ud(w))(U(y)-U(w'))\big\}\notag\\
&-2\,\KK(z-w)\cdot\KK(z-w')\,K(x,y,z)\,\Tr\big\{(\Ud(z)-\Ud(w))(U(z)-U(w'))\big\}\notag\\
&+\KK(z-w)\cdot\KK(y-w')\,\big(2K(x,y,z)-K(x,x,z)\big)\notag\\
&\hspace{6em}\times\Tr\big\{(\Ud(z)-\Ud(w))(U(y)-U(w'))\big\}\notag\\
&+\KK(x-w)\cdot\KK(z-w')\,\big(2K(x,y,z)-K(y,y,z)\big)\notag\\
&\hspace{6em}\times\Tr\big\{(\Ud(x)-\Ud(w))(U(z)-U(w'))\big\}\Big].
\label{eq:LOBFKL}
\end{align}
In Eq.~\eqref{eq:Presult}, $\delta Y\to\delta Y_P=\log(\vee/k^+)$. The second line is the measured gluon emitted by the harder gluon at $z$. It requires $p^+>k^+$ and therefore carries only $\delta Y_P$.

\subsection{JIMWLK}

Write $t$ for the position of the soft gluon and define
\begin{equation}
P\equiv\frac{\alpha_s^2\,\delta Y}{8\pi^6(N_c^2-1)}\int_{x,y,t}\KK(x-w)\cdot\KK(y-w')\,\langle\rho^f(x)\rho^f(y)\rangle,\quad M(u,v)\equiv\big(\Ud(u)-\Ud(t)\big)\big(U(v)-U(t)\big),
\end{equation}
with the signs $\sigma(x)=\sigma(y)=+1$, $\sigma(w)=\sigma(w')=-1$. Carrying out the derivatives of Eq.~\eqref{eq:HJIMWLK} on Eq.~\eqref{eq:LOx} gives $-\delta Y\,H_{\rm JIMWLK}\otimes{\rm LO}=\sum_{n=1}^5dN^{(n)}$, with
\begin{align}
dN^{(1)}&=P\sum_{v\in\{x,w\}}\ \sum_{u\in\{y,w'\}}\sigma(u)\sigma(v)\,K(u,v,t)\,\Tr\big\{T^e\Ud(v)U(u)T^e\,M(u,v)\big\},\\
dN^{(2)}&=P\sum_{u\in\{x,w\}}\ \sum_{v\in\{y,w'\}}\sigma(u)\sigma(v)\,K(u,v,t)\,\Tr\big\{T^e\Ud(u)U(v)T^e\,M(u,v)^T\big\}=dN^{(1)},\\
dN^{(3)}&=-P\sum_{v\in\{y,w'\}}\sigma(v)\,K(v,v,t)\,\Tr\big\{T^e\big(\Ud(x)-\Ud(w)\big)U(v)T^e\,M(v,v)\big\},\\
dN^{(4)}&=-P\sum_{v\in\{x,w\}}\sigma(v)\,K(v,v,t)\,\Tr\big\{T^e\Ud(v)\big(U(y)-U(w')\big)T^e\,M(v,v)\big\},\\
dN^{(5)}&=+P\Big[\sum_{u\in\{y,w'\}}\sigma(u)K(u,u,t)\,\Tr\big\{(\Ud(u)-\Ud(t))U(u)T^d\big\}\,\Tr\big\{(\Ud(x)-\Ud(w))U(u)T^d\big\}\notag\\
&\qquad\ -\sum_{u\in\{x,w\}}\sigma(u)K(u,u,t)\,\Tr\big\{(\Ud(u)-\Ud(t))U(u)T^d\big\}\,\Tr\big\{\Ud(u)(U(y)-U(w'))T^d\big\}\Big].
\end{align}
In Eq.~\eqref{eq:Tresult}, $\delta Y\to\delta Y_T=\log(k^+/\Lambda)$. Two features deserve emphasis. First, the matrix $M$ is evaluated at the same points $u,v$ at which the derivatives act. In particular, it is evaluated at $w$ or $w'$ when the soft gluon couples to the measured gluon. Second, $dN^{(5)}$, which comes from the term linear in $J_R$ that appears when all derivatives are moved to the right of the Wilson lines, enters with a positive sign. The four terms with one derivative on each side of the cut form $dN^{(1)}+dN^{(2)}$. The four terms with both derivatives on the same side form $dN^{(3)}$, $dN^{(4)}$ and $dN^{(5)}$.
